% GENERATED FILE. Edit canonical lessons, then run npm run generate:books.
% canonical-source-hash: fnv1a64:66098753d37b9af6
% canonical-lessons: ES-C721-rasgo, ES-C721-personalidad, ES-C721-temperamento, ES-C721-defecto, ES-C721-valorar

\chapter{Rasgos de personalidad}
\label{ch:es-rasgos-de-personalidad}
\hlchaptermodality{721}

\begin{chapteropening}
\textbf{By the end of this chapter:} \emph{I can describe someone's character, their strong points and their faults.}

Complete the last lesson of chapter 721: I can describe someone's character, their strong points and their faults.
\end{chapteropening}

\section[el rasgo]{el rasgo — the trait, the feature}

\label{lesson:ES-C721-rasgo}



\begin{quote}
Before the new one: say the Spanish for gratitude, then the Spanish for to appreciate.
\end{quote}

\subsection*{You'll want to know: el rasgo}
\textbf{el rasgo} — \textquotedblleft{}the trait, the feature\textquotedblright{}.

\begin{quote}
La paciencia es su mejor rasgo.
\end{quote}

\begin{quote}
\emph{Patience is her best trait.}
\end{quote}

\emph{Los rasgos} are also facial features: \emph{rasgos finos}.

\begin{cousinweb}
\emph{rasgo} is built on \emph{rasgar} \textquotedblleft{}to tear\textquotedblright{}, which is usually traced to Latin \emph{resecāre} \textquotedblleft{}to cut off\textquotedblright{}, from \emph{secāre} \textquotedblleft{}to cut\textquotedblright{}. A \emph{rasgo} was first a stroke of the pen, then a feature.
\end{cousinweb}

\begin{grammarlens}[title={What you've built}]
One more word for telling, explaining and describing at B1.
\end{grammarlens}

\subsection*{Your turn}
\begin{itemize}
\raggedright
  \item \emph{Say it:} \emph{el rasgo}
  \item \emph{Say it:} \emph{el rasgo}, once more
  \item \emph{Say it:} a sentence of your own with \emph{el rasgo}
  \item \emph{Recall:} say the Spanish for gratitude, then the Spanish for to appreciate, then say \emph{el rasgo} again
\end{itemize}

\begin{tcolorbox}[breakable,colback=teal!4,colframe=teal!35!black,title={Before you move on}]
What does el rasgo mean? (\textquotedblleft{}The trait, the feature\textquotedblright{}.) Say it once more.
\end{tcolorbox}
\section[la personalidad]{la personalidad — personality}

\label{lesson:ES-C721-personalidad}



\begin{quote}
Before the new one: say the Spanish for to appreciate, then the Spanish for the trait.
\end{quote}

\subsection*{You'll want to know: la personalidad}
\textbf{la personalidad} — \textquotedblleft{}personality\textquotedblright{}.

\begin{quote}
Tiene una personalidad muy fuerte.
\end{quote}

\begin{quote}
\emph{She has a very strong personality.}
\end{quote}

\begin{cousinweb}
\emph{personalidad} is a learned borrowing from Late Latin \emph{persōnālitās} (EtymDB), built on \emph{persōna} \textquotedblleft{}an actor's mask, a role\textquotedblright{}, the source of \emph{persona}.
\end{cousinweb}

\begin{grammarlens}[title={What you've built}]
One more word for telling, explaining and describing at B1.
\end{grammarlens}

\subsection*{Your turn}
\begin{itemize}
\raggedright
  \item \emph{Say it:} \emph{la personalidad}
  \item \emph{Say it:} \emph{la personalidad}, once more
  \item \emph{Say it:} a sentence of your own with \emph{la personalidad}
  \item \emph{Recall:} say the Spanish for to appreciate, then the Spanish for the trait, then say \emph{la personalidad} again
\end{itemize}

\begin{tcolorbox}[breakable,colback=teal!4,colframe=teal!35!black,title={Before you move on}]
What does la personalidad mean? (\textquotedblleft{}Personality\textquotedblright{}.) Say it once more.
\end{tcolorbox}
\section[el temperamento]{el temperamento — temperament}

\label{lesson:ES-C721-temperamento}



\begin{quote}
Before the new one: say the Spanish for the trait, then the Spanish for personality.
\end{quote}

\subsection*{You'll want to know: el temperamento}
\textbf{el temperamento} — \textquotedblleft{}temperament\textquotedblright{}.

\begin{quote}
Es un niño de temperamento tranquilo.
\end{quote}

\begin{quote}
\emph{He is a child with a calm temperament.}
\end{quote}

\emph{Tener temperamento} — to be fiery.

\begin{cousinweb}
\emph{temperamento} is a learned borrowing from Latin \emph{temperāmentum} \textquotedblleft{}a due mixture\textquotedblright{} (EtymDB), from \emph{temperāre} \textquotedblleft{}to mix in proportion\textquotedblright{}; old medicine saw character as a mix of body fluids.
\end{cousinweb}

\begin{grammarlens}[title={What you've built}]
One more word for telling, explaining and describing at B1.
\end{grammarlens}

\subsection*{Your turn}
\begin{itemize}
\raggedright
  \item \emph{Say it:} \emph{el temperamento}
  \item \emph{Say it:} \emph{el temperamento}, once more
  \item \emph{Say it:} a sentence of your own with \emph{el temperamento}
  \item \emph{Recall:} say the Spanish for the trait, then the Spanish for personality, then say \emph{el temperamento} again
\end{itemize}

\begin{tcolorbox}[breakable,colback=teal!4,colframe=teal!35!black,title={Before you move on}]
What does el temperamento mean? (\textquotedblleft{}Temperament\textquotedblright{}.) Say it once more.
\end{tcolorbox}
\section[el defecto]{el defecto — the fault, the flaw}

\label{lesson:ES-C721-defecto}



\begin{quote}
Before the new one: say the Spanish for personality, then the Spanish for temperament.
\end{quote}

\subsection*{You'll want to know: el defecto}
\textbf{el defecto} — \textquotedblleft{}the fault, the flaw\textquotedblright{}.

\begin{quote}
Su mayor defecto es la impaciencia.
\end{quote}

\begin{quote}
\emph{His biggest fault is impatience.}
\end{quote}

The opposite is \emph{la virtud} or \emph{la cualidad}.

\begin{cousinweb}
\emph{defecto} is a learned borrowing from Latin \emph{dēfectus} \textquotedblleft{}failure, lack\textquotedblright{} (EtymDB), from \emph{dēficere} \textquotedblleft{}to fail\textquotedblright{}, built on \emph{dē-} and \emph{facere} \textquotedblleft{}to do\textquotedblright{}.
\end{cousinweb}

\begin{grammarlens}[title={What you've built}]
One more word for telling, explaining and describing at B1.
\end{grammarlens}

\subsection*{Your turn}
\begin{itemize}
\raggedright
  \item \emph{Say it:} \emph{el defecto}
  \item \emph{Say it:} \emph{el defecto}, once more
  \item \emph{Say it:} a sentence of your own with \emph{el defecto}
  \item \emph{Recall:} say the Spanish for personality, then the Spanish for temperament, then say \emph{el defecto} again
\end{itemize}

\begin{tcolorbox}[breakable,colback=teal!4,colframe=teal!35!black,title={Before you move on}]
What does el defecto mean? (\textquotedblleft{}The fault, the flaw\textquotedblright{}.) Say it once more.
\end{tcolorbox}
\section[valorar]{valorar — to value, to appreciate}

\label{lesson:ES-C721-valorar}



\begin{quote}
Before the new one: say the Spanish for temperament, then the Spanish for the fault.
\end{quote}

\subsection*{You'll want to know: valorar}
\textbf{valorar} — \textquotedblleft{}to value, to appreciate\textquotedblright{}.

\begin{quote}
Valoro mucho a mis amigos.
\end{quote}

\begin{quote}
\emph{I really value my friends.}
\end{quote}

Also to assess: \emph{valorar los daños}.

\begin{cousinweb}
\emph{valorar} is built on \emph{valor} \textquotedblleft{}worth\textquotedblright{}, which comes from Late Latin \emph{valor}, from \emph{valēre} \textquotedblleft{}to be strong, to be worth\textquotedblright{} (EtymDB).
\end{cousinweb}

\begin{grammarlens}[title={What you've built}]
One more word for telling, explaining and describing at B1.
\end{grammarlens}

\subsection*{Your turn}
\begin{itemize}
\raggedright
  \item \emph{Say it:} \emph{valorar}
  \item \emph{Say it:} \emph{valorar}, once more
  \item \emph{Say it:} a sentence of your own with \emph{valorar}
  \item \emph{Recall:} say the Spanish for temperament, then the Spanish for the fault, then say \emph{valorar} again
\end{itemize}

\begin{tcolorbox}[breakable,colback=teal!4,colframe=teal!35!black,title={Before you move on}]
What does valorar mean? (\textquotedblleft{}To value, to appreciate\textquotedblright{}.) Say it once more.
\end{tcolorbox}
